\documentclass[11pt,aspectratio=169]{beamer}
\usepackage[T1]{fontenc}
\usepackage[utf8]{inputenc}
\usepackage[english]{babel}
\usepackage{mathpazo}
\usefonttheme{serif}
\usepackage{amsmath,amssymb,booktabs,array,tabularx,graphicx,microtype}
\usepackage{adjustbox}
\usepackage{tikz}
\usetikzlibrary{arrows.meta,calc,positioning,decorations.pathreplacing,patterns,fit,backgrounds}

% Figures come from the run of analyses/groups; the TikZ figures are in tikz/.
\graphicspath{{../../../analyses/groups/out/d120_eh500_ez300_th45_g120_q0.1/}}

% Colours of the report: S, G and A; neutrals for text, grids and fills.
\definecolor{blueI}{HTML}{2A78D6}
\definecolor{orangeJ}{HTML}{EB6834}
\definecolor{aqua}{HTML}{1BAF7A}
\definecolor{ink}{HTML}{0B0B0B}
\definecolor{ink2}{HTML}{52514E}
\definecolor{muted}{HTML}{9A9893}
\definecolor{light}{HTML}{E6E5E1}
\definecolor{steel}{RGB}{74,96,121}
\definecolor{steellight}{RGB}{192,206,220}

\tikzset{
  >={Stealth[length=1.8mm]},
  every picture/.style={line cap=round, line join=round},
  lbl/.style={font=\footnotesize},
  note/.style={font=\footnotesize, text=ink2},
}

\newcommand{\Sl}{\mathrm{S}}
\newcommand{\G}{\mathrm{G}}
\newcommand{\A}{\mathrm{A}}
\newcommand{\SG}{\Sl\!\mid\!\G}
\newcommand{\SAA}{\Sl\!\mid\!\A}
\newcommand{\Hloc}{H_{\mathrm{loc}}}
\newcommand{\MSD}{\mathrm{MSD}}

% Theme: steel-blue title bar, white pages, as in the earlier decks.
\setbeamercolor{structure}{fg=steel}
\setbeamercolor{normal text}{fg=ink,bg=white}
\setbeamercolor{frametitle}{fg=white,bg=steel}
\setbeamercolor{alerted text}{fg=blueI}
\setbeamercolor{takeaway}{fg=steel,bg=steellight!30}
\setbeamercolor{defbox}{fg=ink,bg=steellight!28}
\setbeamersize{text margin left=7mm,text margin right=7mm}
\setbeamertemplate{navigation symbols}{}
\setbeamertemplate{itemize item}{\textcolor{steel}{\small$\bullet$}}
\setbeamertemplate{itemize subitem}{\textcolor{steel}{\textendash}}
\setbeamertemplate{itemize/enumerate body begin}{\setlength{\itemsep}{0.3em}}
\setbeamertemplate{frametitle}{%
  \nointerlineskip
  \begin{beamercolorbox}[wd=\paperwidth,ht=9mm,dp=3mm,center]{frametitle}
    \large\insertframetitle\strut
  \end{beamercolorbox}
  \begin{tikzpicture}[remember picture,overlay]
    \node[anchor=north east,text=white,inner sep=0pt,xshift=-2.3mm,yshift=-4.2mm]
      at(current page.north east){\small\insertframenumber};
  \end{tikzpicture}}
\setbeamertemplate{footline}{}

\newcommand{\source}[1]{\par\vspace{1mm}{\fontsize{7}{8.3}\selectfont\color{ink2}#1\par}}
\newcommand{\takeaway}[2][\linewidth]{\par\vfill
  \begin{beamercolorbox}[wd=#1,sep=5pt]{takeaway}\small #2\end{beamercolorbox}}
% A definition: its name in bold, then the statement, on a pale band.
\newcommand{\defn}[2]{%
  \begin{beamercolorbox}[wd=\linewidth,sep=5pt]{defbox}\small\textcolor{steel}{\textbf{#1.}}\enspace #2\end{beamercolorbox}}
% A TikZ figure, scaled to fit the given height and the line width.
\newcommand{\tikzfig}[2][0.8\textheight]{%
  \begin{center}\adjustbox{max width=\linewidth,max height=#1}{\input{tikz/#2}}\end{center}}
% A figure of the run at full width, at most the given height.
\newcommand{\fig}[2][0.5\textheight]{%
  \begin{center}\includegraphics[width=\linewidth,height=#1,keepaspectratio]{#2}\end{center}}
% A dense figure on a page of its own, with a small title.
\newcommand{\plainfig}[3][fig]{%
\begin{frame}[plain,label=#1]
\vspace{-1mm}{\small\textcolor{steel}{\textbf{#2}}}\par\vspace{1mm}
\begin{center}\includegraphics[height=0.86\paperheight,width=\linewidth,keepaspectratio]{#3}\end{center}
\end{frame}}
% Small links from a slide of the talk to the appendix slides that detail it.
\newcommand{\applink}[2]{\hyperlink{#1}{\textcolor{blueI}{$\to$ #2}}}
% The same links in one row along the bottom of the page.
\newcommand{\apprefrow}[1]{\begin{tikzpicture}[remember picture,overlay]
  \node[anchor=south east,inner sep=0,font=\fontsize{7}{8.5}\selectfont]at([xshift=-7mm,yshift=2.5mm]current page.south east){#1};
  \end{tikzpicture}}
\newcommand{\appref}[1]{\begin{tikzpicture}[remember picture,overlay]
  \node[anchor=south east,inner sep=0,align=right,font=\fontsize{7}{8.5}\selectfont]at([xshift=-7mm,yshift=2.5mm]current page.south east){#1};
  \end{tikzpicture}}
% A big number with its caption, in a column.
\newcommand{\bignum}[3]{{\centering{\Huge\textcolor{#1}{#2}}\par\smallskip\small #3\par}}
% A full-page divider that opens one part of the talk.
\newcommand{\partframe}[2]{%
\begin{frame}[plain,noframenumbering]
\begin{tikzpicture}[remember picture,overlay]
 \fill[steel](current page.south west)rectangle(current page.north east);
 \node[text=white,align=center,text width=.82\paperwidth]at([yshift=4mm]current page.center)
   {{\Large\bfseries #1}\\[3mm]{\normalsize\color{steellight}#2}};
\end{tikzpicture}
\end{frame}}

\hypersetup{pdftitle={Group time against solo flight},pdfauthor={Matteo Di Sante},
  colorlinks=true,urlcolor=blueI,linkcolor=blueI}

\begin{document}

\begin{frame}[plain]
\begin{tikzpicture}[remember picture,overlay]
 \fill[steel](current page.south west)rectangle(current page.north east);
 \fill[steellight](current page.south west)rectangle([yshift=20mm]current page.south east);
 \node[text=white,align=center,text width=.9\paperwidth]at([yshift=8mm]current page.center)
   {{\LARGE\bfseries Does a paraglider move differently in a group?}\\[4mm]
   {\small Group time against solo flight, on open flights}};
 \node[text=steel,align=center]at([yshift=10mm]current page.south)
   {{\normalsize 3 November 2026}\\[1.5mm]};
\end{tikzpicture}
\end{frame}

%% ------------------------------------------------------------------ main

\begin{frame}{Four questions}
\begin{itemize}\setlength{\itemsep}{1.0em}\large
  \item[\textcolor{steel}{\Huge\textbf{1}}] Do flights in a group move differently from solo flights?\\{\normalsize\textcolor{steel}{Raw}}
  \item[\textcolor{steel}{\Huge\textbf{2}}] With route, season, hour and wing fixed, does the group still make a difference?\\{\normalsize\textcolor{steel}{All fixed}}
  \item[\textcolor{steel}{\Huge\textbf{3}}] How much does each of these factors count?\\{\normalsize\textcolor{steel}{Each factor}}
  \item[\textcolor{steel}{\Huge\textbf{4}}] Does the same flight change when it enters a group, and when it leaves?\\{\normalsize\textcolor{steel}{Same flight}}
\end{itemize}
\end{frame}

\begin{frame}{Three populations of flight time}
\tikzfig[0.72\textheight]{populations}
\begin{itemize}
  \item[\textcolor{blueI}{$\bullet$}] $\Sl$ (Solo flight): flights that never entered a group.
  \item[\textcolor{orangeJ}{$\bullet$}] $\G$ (Group flight): group time of the flights that entered one.
  \item[\textcolor{aqua}{$\bullet$}] $\A$ (Solo stretch): the solo part of the group flights.
\end{itemize}
\takeaway{A \emph{piece} is an uninterrupted stretch of one population.}
\end{frame}

\begin{frame}{Which flights: the open ones}
\vspace{-3mm}
\definecolor{bordeaux}{HTML}{8C2F0F}
\begin{center}\begin{tikzpicture}
  \node[anchor=south west,inner sep=0](img)at(0,0){\includegraphics[width=\linewidth,height=0.57\textheight,keepaspectratio]{open_declared/opening.pdf}};
  \begin{scope}[x={(img.south east)},y={(img.north west)}]
    \draw[bordeaux,thick](0.124,0.92)ellipse(0.04 and 0.11);
  \end{scope}
\end{tikzpicture}\end{center}
\begin{itemize}
  \item \textbf{Opening} $=\dfrac{\lVert\mathbf r_{\mathrm{land}}-\mathbf r_{\mathrm{take}}\rVert}{\sum_i\lVert\mathbf r_{i+1}-\mathbf r_i\rVert}$: 1 straight, $\approx0$ back at the takeoff. \vspace{2mm}
  \item We keep opening $\ge0.1$: 84\% of the flights; 88\% of low takeoffs, 80\% of high ones (right).
\end{itemize}
\end{frame}

%% Question 1
\partframe{Question 1}{Do flights in a group move differently from solo flights?}

\begin{frame}{Fitted $H$ and $\Hloc$ of S, A and G}
\vspace{-5mm}
\fig[0.60\textheight]{open_min0.1/slide_raw_msd_fits.pdf}
\vspace{-3mm}
\begin{itemize}\footnotesize\setlength{\itemsep}{0.2em}
  \item $\mathrm{MSD}\sim\tau^{2\Hloc}$.
  \item Each fit on the pieces of at least end lag / 0.8: \textbf{flights constant along the lags, fewer for longer fits}.
\end{itemize}
\begin{center}\scriptsize
\renewcommand{\arraystretch}{1.1}
\begin{tabular}{@{}l c c c@{}}
\toprule
 & to 1000~s & to 2760~s & to 11\,450~s \\
\midrule
$H_\G-H_\Sl$ & $+0.002\ [+0.001,+0.004]$ & $+0.009\ [+0.007,+0.011]$ & $+0.019\ [+0.006,+0.032]$ \\
\bottomrule
\end{tabular}
\end{center}
\vspace{-1mm}
\end{frame}

\begin{frame}{Raw: group time is straighter at the long lags}
\vspace{-3mm}
\fig[0.41\textheight]{open_min0.1/slide_raw_delta.pdf}
\vspace{-2mm}
\begin{itemize}
  \item $\Hloc(\tau)$: local exponent of the $\mathrm{MSD}\sim\tau^{2\Hloc}$.
  \item $\Hloc(1000~\mathrm{s})$: $\G$ \textcolor{orangeJ}{0.915}, $\Sl$ \textcolor{blueI}{0.897}, $\A$ \textcolor{aqua}{0.894}.
  \item \fcolorbox{steel}{white}{$\G-\Sl=+0.018\ [+0.015,+0.021]$; positive from 170~s}
\end{itemize}
\vspace{2mm}
{\scriptsize\color{ink2} Same flights at every lag: pieces of at least 1250~s, so lags up to 1000~s.\par}
\takeaway{At first the group waits to form, a solo pilot leaves at once. After one glide (median $\approx200$~s) group time gains, probably because group flights favour zones and days with more thermals.}
\vspace{2mm}
\apprefrow{\applink{app-hloc}{how we compute $\Hloc$}\quad\applink{app-bands}{the bands}\quad\applink{app-rawlong}{longer lags, other pools}\quad\applink{app-map}{map of flights}}
\end{frame}

\begin{frame}{Raw: the durations of the phases}
\fig[0.72\textheight]{open_min0.1/raw_survival.pdf}
\begin{itemize}
  \item Climbs, searches and glides last as long in S, A and G: the survival curves overlap.
\end{itemize}
\end{frame}


\begin{frame}{Raw: group time also flies differently}
\vspace{-2mm}
\begin{columns}[T,onlytextwidth]
\column{.67\textwidth}\includegraphics[width=\linewidth,height=0.70\textheight,keepaspectratio]{open_min0.1/slide_raw_phases.pdf}
\column{.30\textwidth}\footnotesize
\begin{itemize}\setlength{\itemsep}{0.25em}
  \item[\textcolor{steel}{\textbf{1, 2}}] $\G-\Sl\approx0.02$, with pieces of at least 1250~s (1) or 3450~s (2).
  \item[\textcolor{steel}{\textbf{3, 6}}] $\G$ climbs more and finds more climbs.
  \item[\textcolor{steel}{\textbf{4, 5}}] $\G$ searches less and better.
  \item[\textcolor{steel}{\textbf{7}}] $\A\approx\G$: zones with stronger thermals.
  \item[\textcolor{steel}{\textbf{8}}] $\G$ keeps a higher glide speed (Williams et al.\ 2018).
  \item[\textcolor{steel}{\textbf{9}}] $\G$ glides flatter.
\end{itemize}
\vspace{0.3em}{\scriptsize\color{ink2} pp: percentage points; 90\% bands. $\G-\Sl$ excludes 0 in panels 3--9.\par}
\end{columns}
\noindent\parbox{.72\textwidth}{\scriptsize\color{ink2} Flights S / A / G: panel 1, 43\,977 / 8\,991 / 11\,077; panel 2, 33\,613 / 6\,235 / 6\,297; panels 3--9 (with phase time), 27\,665 / 6\,397 / 7\,896.}
\appref{\applink{app-statdef}{how they are computed}\\\applink{app-estimator}{how each flight counts}}
\end{frame}



%% Question 2
\partframe{Question 2}{With route, season, hour and wing fixed, does the group still make a difference?}

\begin{frame}{The raw gap mixes everything}
\begin{columns}[T,onlytextwidth]
\Large
\column{.5\textwidth}
Group time differs from solo time in
\begin{itemize}
  \item the \textbf{route} flown,
  \item the \textbf{atmospheric conditions} (the season and the hour),
  \item the \textbf{wing},
  \item the pilots who fly in groups.
\end{itemize}
\column{.46\textwidth}
We take them apart
\begin{enumerate}
  \item fix them, and see what is left;
  \item split the gap among them.
\end{enumerate}
\end{columns}
\takeaway{\Large\strut Same route, season, hour and wing: does the group still make a difference?\strut}
\end{frame}

\begin{frame}{The controls $\SG$ and $\SAA$: same route, season, hour and wing}
\vspace{-12mm}
\tikzfig[0.62\textheight]{matching}
\begin{itemize}
  \item $\SG$ answers: \textbf{how does a solo pilot fly the stretch that $\G$ flew?} For each piece of $\G$, the segments of solo flights along the same route, same direction. $\SAA$: the same for a piece of $\A$.
  \item A segment is kept only if it meets the four conditions on the right.
  \item Tolerances chosen to keep at least 200 pieces of $\G$ per wing level (beg/exp); $\A$ keeps 123 beginners on long pieces.
\end{itemize}
\apprefrow{\applink{app-matching}{what is a control}\quad\applink{app-longT}{why the long dataset stops at 2760~s}}
\end{frame}

\begin{frame}{On the same route, group pilots are straighter, in a group or alone}
\vspace{-1mm}
\fig[0.56\textheight]{open_min0.1/slide_matched_deltas.pdf}
\vspace{-7mm}
\begin{center}\scriptsize\renewcommand{\arraystretch}{0.95}
\begin{tabular}{@{}l l l@{}}
\toprule
90\% bands & $\Hloc(1000~\mathrm{s})$ & $\Hloc(1840~\mathrm{s})$ \\
\midrule
$\G-\SG$ = the group + pilots and day & $+0.010\ [+0.004,+0.017]$ & $+0.021\ [+0.008,+0.034]$ \\
$\A-\SAA$ = the same pilots and days, without the group & $+0.006\ [-0.003,+0.016]$ & $+0.027\ [+0.008,+0.046]$ \\
\textbf{difference = the group alone, $\boldsymbol{\approx0}$} & $\boldsymbol{+0.004\ [-0.007,+0.015]}$ & $\boldsymbol{-0.006\ [-0.028,+0.016]}$ \\
\bottomrule
\end{tabular}\end{center}
{\scriptsize\color{ink2} Pieces of $\G$ with a control: 2\,717 (from 2\,474 flights, 1\,378 pilots) at $\geq1250$~s, 700 at $\geq3450$~s; of $\A$: 1\,607 (1\,504 flights, 990 pilots) and 338. Controls: 8\,094 for $\G$, 4\,522 for $\A$.\par}
\apprefrow{\applink{app-selection}{why the difference}\quad\applink{app-matchedhloc}{$\Hloc$ curves}\quad\applink{app-fittedroute}{fitted $H$}\quad\applink{app-robust}{robustness}}
\end{frame}

\begin{frame}{Same route: net of pilots and day, only the time in climb remains}
\vspace{-3mm}
\begin{columns}[c,onlytextwidth]
\column{.64\textwidth}\includegraphics[width=\linewidth,height=0.76\textheight,keepaspectratio]{open_min0.1/slide_matched_forest.pdf}\par\vspace{1mm}
{\footnotesize\color{ink2} Pieces of $\G$ / $\A$: 2\,717 / 1\,607 for $\Hloc(1000~\mathrm{s})$, 700 / 338 for $\Hloc(1840~\mathrm{s})$, 1\,422 / 831 for the phase statistics.\par}
\column{.34\textwidth}\footnotesize
\begin{itemize}\setlength{\itemsep}{0.6em}
  \item Panels 1--2: at 1840~s $\A-\SAA\approx\G-\SG$. The gap comes with the pilots and the day, not with the group; to fix them, the same flight (question 4).
  \item Only the time in climb survives in the difference: $+2.4$ pp $[+0.8,+4.0]$.
  \item Glide speed: raw $+2.1$ km/h, here $-0.9$ ($\G-\SG$) and $-1.0$ ($\A-\SAA$). \textbf{Probably group flights are on routes where everyone glides faster.}
\end{itemize}
{\scriptsize\color{ink2} 90\% bands.\par}
\end{columns}
\apprefrow{\applink{app-forest}{by wing level}\quad\applink{app-robust}{robustness}}
\end{frame}

%% Question 3
\partframe{Question 3}{How much does each of these factors count?}

\begin{frame}{The eight sets of controls: the route splits them in two}
\vspace{-4mm}
\fig[0.53\textheight]{open_min0.1/decomposition_2760_delta.pdf}
\vspace{-3mm}
\begin{columns}[T,onlytextwidth]
\column{.49\textwidth}\footnotesize
\begin{itemize}
  \item Right: $\Delta\Hloc(\tau)$ of $\G$ minus eight sets of controls, one per subset of route $R$, season and hour $T$, wing $W$ fixed, from $\emptyset$ to $R{+}T{+}W=\SG$.
\end{itemize}
\column{.49\textwidth}\footnotesize
\begin{itemize}
  \item With $R$ (solid): positive at long lags. Without $R$: negative. $T$, $W$: little.
\end{itemize}
\end{columns}
\source{$\emptyset$ is not the raw gap: raw $+0.017$, here $-0.039$ at 1840~s, because the controls are windows of very long solo flights, which are straighter. Bands for $\emptyset$ and $R{+}T{+}W$. Left: $\Hloc$ of $\G$, of the controls of $\emptyset$ and of $\SG$. 700 pieces of $\G$ of at least 3450~s, the same in every set.}
\apprefrow{\applink{app-factors}{the nine quantities, three sets}}
\end{frame}

\begin{frame}{How much each factor explains: the Shapley split}
\vspace{-2mm}
\begin{columns}[T,onlytextwidth]
\column{.60\textwidth}\includegraphics[width=\linewidth,height=0.74\textheight,keepaspectratio]{open_min0.1/slide_shapley.pdf}\par\vspace{1mm}
{\footnotesize\color{ink2} Pieces of $\G$, the same in every set: 2\,717 for $\Hloc(1000~\mathrm{s})$, 700 for $\Hloc(1840~\mathrm{s})$, 1\,422 for the phase statistics.\par}
\column{.38\textwidth}\footnotesize
\begin{itemize}\setlength{\itemsep}{0.6em}
  \item \textbf{Question:} how much of the move from $\Delta(\emptyset)$ to $\Delta(F)=\G-\SG$ comes from each of $R$, $T$, $W$?
  \item \textbf{Method:} the gap is computed for each combination of fixed factors; we average the change from adding $k$ over the 6 orders. \textbf{$\varphi_k$ = minus that average.}
  \item \textbf{They add up:} $\Delta(\emptyset)=\Delta(F)+\varphi_R+\varphi_T+\varphi_W$. Time in climb: $-9.5=+3.0-12.0+0.2-0.7$.
  \item \textbf{The route carries most of the move:} 98\% at $\Hloc(1840~\mathrm{s})$, 96\% time in climb, 84\% glide speed.
\end{itemize}
\end{columns}
\apprefrow{\applink{app-phi}{how $\varphi$ is computed}\quad\applink{app-phisum}{why they add up}}
\end{frame}

%% Question 4
\partframe{Question 4}{Does the same flight change when it enters a group, and when it leaves?}

\begin{frame}{The same flight, before, during and after}
\tikzfig[0.72\textheight]{windows}
\[
  D_\text{in}=\text{start}-\text{pre},\quad D_\text{out}=\text{end}-\text{post},\quad
  g=\tfrac12(D_\text{in}+D_\text{out}),\quad d\,w=\tfrac12(D_\text{in}-D_\text{out})
\]
\begin{itemize}
  \item Same pilot, glider, logger and day: it removes who flies and when$^*$.
  \item $g$: what the group does; $d\,w$: the drift of the flight over time.
\end{itemize}
\source{$^*$Only in part: the windows can be hours apart and over different terrain, so the weather may change.}
\end{frame}

\begin{frame}{Fitted $H$ in the same flight}
\defn{$\Delta H$ of $g$}{$\tfrac12\big[(H_\text{start}-H_\text{pre})+(H_\text{end}-H_\text{post})\big]$}
\fig[0.50\textheight]{open_min0.1/slide_fitted_h_flight.pdf}
\takeaway{$\Delta H$ of $g$ is positive for every window (fit from 10~s to 0.8 of the window).}
\end{frame}

\begin{frame}{With windows of 2400~s the flight is straighter in the group}
\vspace{-3mm}
\fig[0.46\textheight]{open_min0.1/slide_prepost_delta.pdf}
\begin{itemize}
  \item $g=+0.026\ [+0.020,+0.032]$ at $\Hloc(1000~\mathrm{s})$.
  \item Time in climb $+2.8\ [+2.0,+3.6]$ pp.
  \item $g>0$ from 50~s, as $\G-\SG$ on long pieces (60~s); raw $\G-\Sl$: 170~s.
  \item The controls: how much does $\Hloc$ change between two consecutive windows, with no change of state?
\end{itemize}
\source{1\,595 events from 1\,542 flights. Control curves: about 5\,300 pairs in the solo time of group flights, 26\,600 in solo flights, one random pair of consecutive windows per solo stretch of at least 4800~s; not matched to the events.}
\appref{\applink{app-prepostforest}{all the differences}}
\end{frame}

\begin{frame}{The same flight: all the quantities}
\vspace{-3mm}
\begin{columns}[c,onlytextwidth]
\column{.62\textwidth}\includegraphics[width=\linewidth,height=0.9\textheight,keepaspectratio]{open_min0.1/slide_prepost_forest.pdf}
\column{.36\textwidth}\footnotesize
\begin{itemize}\setlength{\itemsep}{0.7em}
  \item \textbf{$g$ excludes 0 for eight of nine quantities.} (Not for $\Hloc(1840~\mathrm{s})$).
  \item The controls are near 0 for $\Hloc$: time alone does not produce $g$.
  \item For climbing, time alone goes the other way.
\end{itemize}
{\scriptsize\color{ink2} Windows of 2400~s; 90\% bands. Events: 1\,595 (1\,542 flights) for $\Hloc$, 1\,150 (1\,113 flights) for the phase statistics.\par}
\end{columns}
\apprefrow{\applink{app-prepostforest}{also $d\,w$ and post $-$ pre}}
\end{frame}

\begin{frame}{$g$ appears at every window length; $\Hloc(1000~\mathrm{s})$ is a fair quote from 1800~s}
\vspace{-3mm}
\fig[0.58\textheight]{open_min0.1/prepost_windows.pdf}
\begin{itemize}
  \item \textbf{$g>0$ for every window}. Plateau for Window $> 1800 s$ .
\end{itemize}
\source{Events with the four windows: 3\,549, 2\,567, 1\,595 and 695.}
\end{frame}

%% Take-home
\begin{frame}{Take-home}
\renewcommand{\arraystretch}{1.1}
\begin{tabularx}{\linewidth}{@{}c X r@{}}
\textcolor{steel}{\LARGE\textbf{1, 4}} & \large A group effect: in the raw data and within one flight
  & \shortstack[r]{{\LARGE\textcolor{blueI}{$+0.018$}\;\textcolor{aqua}{$+0.026$}}\\[1mm]{\scriptsize $\Hloc(1000~\mathrm{s})$: raw, same flight}} \\[4mm]
\textcolor{steel}{\LARGE\textbf{2}} & \large On the same route, group and solo fly alike
  & \shortstack[r]{{\LARGE\textcolor{orangeJ}{$\approx0$}}\\[1mm]{\scriptsize the group alone}} \\[4mm]
\textcolor{steel}{\LARGE\textbf{3}} & \large The route carries the gap
  & \shortstack[r]{{\LARGE\textcolor{steel}{96--98\%}}\\[1mm]{\scriptsize of what the factors explain}} \\[4mm]
\textcolor{steel}{\LARGE\textbf{+}} & \large More time in climb, in every design
  & \shortstack[r]{{\LARGE\textcolor{steel}{$+3$ pp}}\\[1mm]{\scriptsize raw, same route, same flight}} \\
\end{tabularx}
\takeaway{\large\textbf{Reading:} the group helps to find a better path; on the same path it adds little.}
\end{frame}

%% ------------------------------------------------------------------ appendix

\partframe{Appendix}

\begin{frame}[label=app-hloc]{How we compute $\Hloc$ in practice}
\tikzfig[0.40\textheight]{straightness}
\begin{itemize}
  \item Per flight, the time-averaged MSD on the 10-s grid of its pieces; the ensemble MSD is the mean over flights, one weight per flight.
  \item $\MSD(\tau)\sim\tau^{2H}$ on a geometric grid of lags; $\Hloc(\tau)=\tfrac12\times$ least-squares slope of $\log\MSD$ against $\log\tau$ over the lags within $\pm0.25$ decades of $\tau$ (from $\tau/1.78$ to $1.78\,\tau$); with fewer than 3 lags, the 3 nearest.
  \item $H$: the same fit over all lags, from 10~s to the end of the dataset.
\end{itemize}
\end{frame}

\begin{frame}[label=app-hlocedge]{$\Hloc$ at the edge of the lags: a one-sided estimate}
\tikzfig[0.32\textheight]{edge}
\begin{itemize}
  \item At the longest lags $T$ of a dataset the window is cut on its upper side: the slope uses only lags below $T$, so it is centred at a shorter lag than $\tau$ and is asymmetric.
  \item Example: $\Hloc(1840~\mathrm{s})$ on the dataset to $T=2760$~s uses 1035 to 2760~s instead of 1035 to 3272~s; $\Hloc(1000~\mathrm{s})$ on the dataset to 1000~s uses 562 to 1000~s, but 562 to 1778~s on the dataset to 2760~s: complete.
  \item All curves of a dataset share $T$, hence the same cut: differences between populations are like for like, and every bootstrap draw is cut the same way, so the bands include it.
\end{itemize}
\end{frame}


\begin{frame}[label=app-estimator]{The phase statistics: how each flight counts}
Example, glide speed = distance of glides / time in glides, for two flights:
\begin{center}\small
\begin{tabular}{@{}c r r r r@{}}
\toprule
flight & glide distance & glide time & flight time & glide speed \\
\midrule
1 & 20~km & 0.5~h & 1~h & 40~km/h \\
2 & 60~km & 2~h & 4~h & 30~km/h \\
\bottomrule
\end{tabular}
\end{center}
\begin{columns}[T,onlytextwidth]
\column{.48\textwidth}
\defn{Ratio of sums}{$\dfrac{20+60}{0.5+2}=32$~km/h. Flight 2 weighs four times flight 1: it glides four times longer. The longest flights dominate.}
\column{.48\textwidth}
\defn{Ours: average over flights, then ratio}{$\dfrac{\mathrm{mean}(20/1,\;60/4)}{\mathrm{mean}(0.5/1,\;2/4)}=\dfrac{17.5}{0.5}=35$~km/h. A flight weighs by its share of time gliding, not by its length.}
\end{columns}
\vspace{1mm}
\begin{itemize}
  \item With a unit $u$: $r_u=n_u/T_u$, $q_u=d_u/T_u$, stat $=\sum_u w_u r_u\big/\sum_u w_u q_u$; $n_u$, $d_u$ summed over its runs, $T_u$ its flight time.
  \item The MSD also gives one weight per flight. On the data: time in climb of $\Sl$ 32.9\% (ours), 34.2\% (ratio of sums).
\end{itemize}
\end{frame}

\begin{frame}[label=app-statdef]{The seven phase statistics: formulas}
Per flight, with $T$ its flight time and sums over its runs of the phase; then averaged over flights as on the previous page.
\begin{center}\small
\renewcommand{\arraystretch}{1.35}
\begin{tabular}{@{}l l l@{}}
\toprule
statistic & formula & in words \\
\midrule
time in climb (pp) & $100\,T_{\mathrm{climb}}/T$ & share of the flight spent climbing \\
time in search (pp) & $100\,T_{\mathrm{search}}/T$ & share spent looking for the next climb \\
searches $\to$ climb (pp) & $100\,N_{\mathrm{found}}/N_{\mathrm{search}}$ & share of searches directly followed by a climb \\
climbs per hour & $N_{\mathrm{climb}}/T$ & climbs started per hour of flight \\
climb rate (m/s) & $\sum\Delta z_{\mathrm{climb}}\,/\,T_{\mathrm{climb}}$ & height gained over time climbing \\
glide speed (km/h) & $3.6\sum d_{\mathrm{glide}}\,/\,T_{\mathrm{glide}}$ & distance over time gliding \\
height lost (m/km) & $1000\sum h_{\mathrm{lost}}\,/\sum d_{\mathrm{glide}}$ & height lost per km of glide \\
\bottomrule
\end{tabular}
\end{center}
{\small $d_{\mathrm{glide}}$: straight distance from the start to the end of a glide (the chord); $h_{\mathrm{lost}}$: height lost during it; $N_{\mathrm{found}}$: searches whose next phase is a climb.}
\end{frame}

\begin{frame}[label=app-datasets]{Datasets and the $\tau/0.8$ rule}
\begin{columns}[T,onlytextwidth]
\column{.5\textwidth}
\fig[0.72\textheight]{open_min0.1/raw_durations.pdf}
\column{.46\textwidth}
A piece of $\ell$~s gives $(\ell-\tau)/10+1$ pairs of origins. Requiring $\ell\ge\tau/0.8$ gives the same share of origins at the longest lag.
\begin{itemize}
  \item Dataset to $T$: pieces of at least $T/0.8$~s, at every lag up to $T$.
  \item 1000: pieces of 1250~s; 2760: 3450~s; longest (11\,450~s): 14\,312~s.
  \item Pieces of $\G$ are shorter: a group lasts at least $\tau^*=38.1$~min.
\end{itemize}
\end{columns}
\end{frame}

\begin{frame}[label=app-longT]{Why the long dataset stops at 2760~s}
\begin{columns}[T,onlytextwidth]
\column{.56\textwidth}
\footnotesize\setlength{\tabcolsep}{3pt}
\begin{tabular}{@{}l r r r@{}}
\toprule
caliper (month, hour, wing) & $T=2250$ & 2760 & 3380 \\
\midrule
$\pm1$, $\pm1$, class & 311/395 & 166/231 & 74/111 \\
$\pm1$, $\pm1$, level & 327/656 & 176/372 & 79/189 \\
$\pm2$, $\pm2$, class & 513/725 & \textbf{279/421} & 129/228 \\
$\pm2$, $\pm2$, level & 537/1095 & 294/650 & 139/370 \\
\bottomrule
\end{tabular}

{\scriptsize pieces of $\G$ with a control, beginners/experienced; the rule asks for 200 in each level.}
\column{.4\textwidth}
\begin{itemize}
  \item $\Hloc(1840~\mathrm{s})$ uses the lags within $\pm0.25$ decades: 1035 to 3272~s.
  \item At $T=2760$ the window is cut at 2760~s; at 3380~s no caliper reaches 200 pieces per level.
\end{itemize}
\end{columns}
\end{frame}

\begin{frame}[label=app-selection]{Why the difference $(\G-\SG)-(\A-\SAA)$}
\noindent Goal: separate the effect of the group from everything else group flights carry with them: the pilots and the day (its weather beyond season and hour).
\vspace{1mm}
\defn{$\G-\SG$}{Group time against solo controls on the same route, season, hour and wing. It holds the group, and what group flights carry with them: the pilots and the day.}
\vspace{1mm}
\defn{$\A-\SAA$}{The solo stretches of the same flights (same pilots, wings, days) against their controls. There is no group in $\A$: what is left is the pilots and the day.}
\vspace{1mm}
\defn{Difference}{Subtract the second gap from the first: $\underbrace{\G-\SG}_{\text{group}\,+\,\text{pilots, day}}-\underbrace{(\A-\SAA)}_{\text{pilots, day}}=\text{group}$.}
\begin{itemize}
  \item At 1840~s: $+0.021$ for $\G-\SG$, $+0.027$ for $\A-\SAA$, $-0.006$ for the difference: the gap is carried by the pilots and the day.
  \item It assumes pilots and day weigh the same in group time and alone, within the same flight.
\end{itemize}
\end{frame}

\begin{frame}[label=app-phi]{How $\varphi$ is computed}
\vspace{-3mm}
\noindent\small Time in climb (pp). One gap $\Delta(C)=\G-{}$controls for each subset $C$ of $\{R,T,W\}$ fixed; \textbf{fixing $R$ moves it by an amount that depends on what is already fixed.}
\begin{columns}[T,onlytextwidth]
\column{.28\textwidth}
\begin{center}\footnotesize
\renewcommand{\arraystretch}{1.0}
\begin{tabular}{@{}l r@{}}
\toprule
fixed $C$ & $\Delta(C)$ \\
\midrule
$\emptyset$ & $-9.55$ \\
$R$ & $+2.20$ \\
$T$ & $-10.35$ \\
$W$ & $-8.66$ \\
$R{+}T$ & $+2.65$ \\
$R{+}W$ & $+2.64$ \\
$T{+}W$ & $-9.29$ \\
$R{+}T{+}W$ & $+2.96$ \\
\bottomrule
\end{tabular}
\end{center}
\column{.68\textwidth}
\begin{center}\footnotesize
\renewcommand{\arraystretch}{1.0}
\begin{tabular}{@{}l c r c@{}}
\toprule
$R$ added after & gap before $\to$ after & change & weight \\
\midrule
nothing & $-9.55\to+2.20$ & $+11.75$ & $1/3$ \\
$T$ & $-10.35\to+2.65$ & $+13.00$ & $1/6$ \\
$W$ & $-8.66\to+2.64$ & $+11.30$ & $1/6$ \\
$T,W$ & $-9.29\to+2.96$ & $+12.25$ & $1/3$ \\
\bottomrule
\end{tabular}
\end{center}
{\footnotesize Weight: how many of the six orders of fixing $R,T,W$ put $R$ in that state (2, 1, 1, 2 of 6).}
\end{columns}
\begin{itemize}\small
  \item $\varphi_R=-\big(\tfrac13\,11.75+\tfrac16\,13.00+\tfrac16\,11.30+\tfrac13\,12.25\big)=-12.05$.
  \item Same for $T$ ($+0.19$) and $W$ ($-0.65$), from the same eight gaps.
\end{itemize}
\end{frame}

\begin{frame}[label=app-phisum]{Why the three $\varphi$ add up to the whole move}
\vspace{-3mm}
\noindent\small Fix the factors in one order: the three steps add up to $\Delta(F)-\Delta(\emptyset)$, whatever the order. Gain of each step, time in climb (pp):
\begin{center}\footnotesize
\renewcommand{\arraystretch}{1.0}
\begin{tabular}{@{}l r r r r@{}}
\toprule
order & gain of $R$ & gain of $T$ & gain of $W$ & total \\
\midrule
$R,T,W$ & $+11.75$ & $+0.44$ & $+0.31$ & $+12.51$ \\
$R,W,T$ & $+11.75$ & $+0.32$ & $+0.44$ & $+12.51$ \\
$T,R,W$ & $+13.00$ & $-0.80$ & $+0.31$ & $+12.51$ \\
$T,W,R$ & $+12.25$ & $-0.80$ & $+1.06$ & $+12.51$ \\
$W,R,T$ & $+11.30$ & $+0.32$ & $+0.89$ & $+12.51$ \\
$W,T,R$ & $+12.25$ & $-0.63$ & $+0.89$ & $+12.51$ \\
\midrule
mean & $+12.05$ & $-0.19$ & $+0.65$ & $+12.51$ \\
$\varphi=-\text{mean}$ & $-12.05$ & $+0.19$ & $-0.65$ & $-12.51$ \\
\bottomrule
\end{tabular}
\end{center}
\begin{itemize}\small
  \item Each order splits the same total among the three factors; the order changes only the split.
  \item $\varphi_k$: the mean over the six orders of the gain of $k$, opposite sign. The three means add up to the mean of six equal totals.
  \item So $\Delta(\emptyset)=\Delta(F)+\varphi_R+\varphi_T+\varphi_W$ exactly: $-9.5=+3.0-12.05+0.19-0.65$.
\end{itemize}
\end{frame}

\begin{frame}[label=app-factors]{Same route, then all fixed: the nine quantities}
\begin{columns}[c,onlytextwidth]
\column{.64\textwidth}\includegraphics[width=\linewidth,height=0.78\textheight,keepaspectratio]{open_min0.1/slide_factors.pdf}
\column{.34\textwidth}\footnotesize
\begin{itemize}\setlength{\itemsep}{0.6em}
  \item $\G$ against three sets of solo controls: 
  \begin{enumerate}
    \item $\emptyset$, no factor fixed; 
    \item $R$, same route; 
    \item $R{+}T{+}W$, same route, season, hour and wing ($=\SG$).
  \end{enumerate}
  \item Fixing the route moves the gap, often across 0; \textbf{fixing season, hour and wing on top changes little.}
\end{itemize}
\end{columns}
\vfill
{\scriptsize\color{ink2} 90\% bands. ``No factor'': windows of long solo flights, not the raw $\Sl$ of slide 5.\par}
\end{frame}

\begin{frame}[label=app-matching]{Matching: what is a control}
\begin{itemize}
  \item Candidates: open flights with no contact of any length, on another date.
  \item Coverage: at least 80\% of the points of the piece, one every 30~s, have a pool track point within 500~m, with the same direction.
  \item Calipers, tried in order: month $\pm1$ or $\pm2$, hour $\pm1$ or $\pm2$, wing class or level; the chosen one is the first with 200 pieces of $\G$ per level.
  \item Each unit is reweighted to the piece, so a piece with many controls counts once.
\end{itemize}
\source{Isolation (nobody within 10~km) is off in the main result and is one of the robustness lines.}
\end{frame}

\plainfig[app-matchedhloc]{Same route, pieces of at least 1250~s: $\Hloc$ of $\G$, $\SG$, $\A$, $\SAA$ and their differences}{open_min0.1/matched_1000_delta.pdf}

\plainfig[app-matchedlong]{Same route, pieces of at least 3450~s: the same; at 1000~s $\G-\SG=+0.024\ [+0.014,+0.034]$}{open_min0.1/matched_2760_delta.pdf}

\begin{frame}[label=app-fittedroute]{Fitted $H$ on the same route}
\fig[0.56\textheight]{open_min0.1/slide_fitted_h_route.pdf}
\takeaway{$\G-\SG$: $-0.005\ [-0.008,-0.002]$ to 1000~s and $+0.008\ [+0.004,+0.012]$ to 2760~s. The sign depends on the lags: $H$ mixes the small deficit at short lags and the gain at long lags. $\A-\SAA$: about 0 to 1000~s, $+0.006$ to 2760~s.}
\end{frame}

\plainfig[app-forest]{Same route: the nine quantities, by wing level}{open_min0.1/matched_forest.pdf}

\plainfig[app-robust]{Same route, 8 lines: $\Delta\Hloc(1000~\mathrm{s})$ $+0.007$ to $+0.012$; $\Delta\Hloc(1840~\mathrm{s})$ $+0.017$ to $+0.022$ ($+0.031$ to $+0.034$ with a shorter long series)}{open_min0.1/matched_robustness.pdf}

\plainfig[app-map]{Where solo flights and group time are: pieces per 1000~km$^2$}{maps.pdf}

\begin{frame}[label=app-bands]{Where the uncertainty comes from}
\tikzfig[0.24\textheight]{bootstrap}
\begin{itemize}
  \item The archive is one sample: other seasons would bring other pilots and other days, and every number would change.
  \item Flights are not independent: those of one pilot share skill, wing and habits; those of one block share the day and the group.
  \item Pilots and blocks cross: a pilot flies in many blocks, a block holds many pilots.
\end{itemize}
\defn{Block}{flights tied by an \emph{episode} of the group identification (groups linked in time through a shared member), or by a \emph{site-day} (takeoffs within 3~km on one day); chained into connected components.}
\source{A cell is one pilot in one block, about one flight.}
\end{frame}

\begin{frame}[label=app-bandcalc]{The band: three bootstraps, one formula}
\begin{itemize}
  \item Draw the pilots with replacement; each unit weighs the count of its pilot; recompute the statistic. Its variance over the draws is $V_P$.
  \item Same for the blocks ($V_B$) and for the cells ($V_C$): 1\,000 draws each.
\end{itemize}
\defn{Variance and band}{units of one pilot or one block covary, so
\[
  V=\max\bigl(V_P+V_B-V_C,\ V_P,\ V_B\bigr),\qquad \text{band}=\hat\theta\pm1.645\sqrt V\quad(90\%).
\]
$V_C$ is subtracted once: the pairs of one cell are in both $V_P$ and $V_B$.}
\vspace{1mm}
\begin{itemize}
  \item All the numbers of a section use the same draws: a difference has the band of its own draws.
  \item Pieces and their controls are drawn together, control flights included: a control that serves many pieces moves them all at once.
\end{itemize}
\source{Cameron, Gelbach and Miller, J. Bus. Econ. Stat. 29, 238 (2011).}
\end{frame}

\begin{frame}[label=app-coverage]{Does the band hold? A check by simulation}
\begin{columns}[T,onlytextwidth]
\column{.4\textwidth}
\begin{itemize}
  \item On the real designs, replace every value by pilot $+$ block $+$ flight $+$ unit effects: no group effect.
  \item A 90\% band must hold 0 in 90\% of the replicates.
  \item The two-way band does so in 88--99\%; the bands used before did not (59\% when pilots matter).
\end{itemize}
\column{.56\textwidth}
\scriptsize\setlength{\tabcolsep}{2.5pt}
\begin{tabular}{@{}l r r r r r r@{}}
\toprule
 & \multicolumn{2}{c}{raw} & \multicolumn{2}{c}{same route} & \multicolumn{2}{c}{same flight} \\
 & new & old & new & old & new & old \\
\midrule
unit only & 91 & 90 & 99 & 100 & 90 & 88 \\
flight only & 92 & 91 & 89 & 100 & 90 & 90 \\
half flight, half unit & 88 & 85 & 92 & 100 & 88 & 88 \\
pilot strong & 88 & \textcolor{orangeJ}{59} & 88 & 100 & 89 & 81 \\
block strong & 88 & 84 & 88 & 99 & 89 & 88 \\
\bottomrule
\end{tabular}

{\scriptsize \% of replicates whose 90\% band holds 0.}
\end{columns}
\takeaway{Left out of the band: the rules of the matching (varied in the robustness lines) and the dependence on one year (planned).}
\end{frame}

\plainfig[app-prepostforest]{Before and after, windows of 2400~s: all the differences}{open_min0.1/prepost_2400_forest.pdf}



\begin{frame}[label=app-poolvar]{The same plots on a pool that shrinks with the lag}
\vspace{-3mm}
\fig[0.52\textheight]{open_min0.1/slide_pool_variable.pdf}
\begin{itemize}
  \item At each lag the pieces of at least lag / 0.8: the flights fall from 44\,774 / 10\,712 / 11\,105 at 10~s to 4\,309 / 602 / 115 at 11\,450~s (S / A / G).
  \item Beyond about 5000~s $\Hloc$ exceeds 1: the curve follows who is left, the longest flights, not the lag. This is why the main plots use the same pieces at every lag.
  \item $\G-\Sl$ turns positive at 210~s, as in the main plots; the bands are very wide at the longest lags.
\end{itemize}
\end{frame}

\begin{frame}[label=app-poolshort]{The same plots on short pieces only}
\vspace{-3mm}
\fig[0.52\textheight]{open_min0.1/slide_pool_short.pdf}
\begin{itemize}
  \item Pieces of 1000 to 1250~s, the same at every lag up to 800~s: 763 flights of $\Sl$, 977 of $\A$, 80 of $\G$.
  \item Here $\G-\Sl$ is negative, $-0.01$ to $-0.02$ up to about 300~s, and goes to 0 at 800~s: no gain of $\G$ on short pieces. With 80 flights the bands of $\G$ are wide.
\end{itemize}
\end{frame}

\begin{frame}[label=app-rawlong]{Raw at longer lags: dataset to 2760~s}
\fig[0.72\textheight]{open_min0.1/slide_raw_2760_delta.pdf}
\source{Same flights at every lag: pieces of at least 3450~s, so lags up to 2760~s.}
\apprefrow{\applink{app-poolvar}{shrinking pool}\quad\applink{app-poolshort}{short pieces}}
\end{frame}

\begin{frame}[label=app-rawlongest]{Raw at the longest lags: dataset to 11\,450~s}
\fig[0.72\textheight]{open_min0.1/slide_raw_longest_delta.pdf}
\begin{itemize}
  \item $\Hloc$ of $\Sl$ and $\A$ falls at the longest lags, that of $\G$ does not.
  \item $\G$ keeps few flights there: pieces of at least 14\,312~s, the same at every lag up to 11\,450~s.
\end{itemize}
\end{frame}

\end{document}
